\NSPage{025}%
determine a unique \NSi{NS-F-P025-001}{q>|z|^{1/D}}. Indeed,
\NSi{NS-F-P025-002}{q-z^2q^{2h}} is zero at that lower endpoint, tends to
infinity, and has derivative
\NSi{NS-F-P025-003}{1-2hz^2q^{2h-1}=L\geq 1-2h>0} on the indicated interval.
Thus \NSi{NS-F-P025-004}{|\eta|<1}; the parameter endpoints
\NSi{NS-F-P025-005}{\eta=\pm1} describe one-sided limits. Throughout the
profile construction, we write
\NSi{NS-F-P025-006}{\mathcal D_Xf=X\partial_Xf} for differentiation in the logarithmic
radial coordinate, with \NSi{NS-F-P025-007}{\eta} held fixed. To compute the
residual of the field described in Section~\ref{sec:3.1}, we now derive the
operators converting time and axial derivatives into derivatives in
\NSi{NS-F-P025-008}{(X,\eta)}.

\begin{lemma}\label{lem:4.1}
For a smooth profile \NSi{NS-F-P025-009}{f} and any real
\NSi{NS-F-P025-010}{b},
\NSFormula{NS-F-P025-011}{display}{4.2}
\begin{equation}
\begin{gathered}
 \partial_t(q^bf)=q^{b-1}T_bf,\qquad
 \partial_z(q^bf)=q^{b-D}Z_bf,\qquad
 \left\{
 \begin{aligned}
 T_bf&=L^{-1}(-bf+D\eta f_\eta+\mathcal D_Xf),\\
 Z_bf&=L^{-1}(2b\eta f+df_\eta-2\eta \mathcal D_Xf).
 \end{aligned}
 \right.
\end{gathered}
\tag{4.2}\label{eq:4.2}
\end{equation}
\end{lemma}

\begin{proof}
Differentiating the equation for \NSi{NS-F-P025-012}{q}, followed by
\NSi{NS-F-P025-013}{\eta=zq^{-D}} and
\NSi{NS-F-P025-014}{X=s/q}, gives
\NSFormula{NS-F-P025-015}{display}{none}
\[
\begin{aligned}
 q_t&=-L^{-1},& \eta_t&=D\eta/(qL),& X_t&=X/(qL),\\
 q_z&=2\eta q^{1-D}/L,& \eta_z&=d/(q^DL),& X_z&=-2\eta X/(q^DL).
\end{aligned}
\]
The middle entry in the second row uses
\NSi{NS-F-P025-016}{L-2D\eta^2=d}. The chain rule proves the claim.
\end{proof}

To construct the leading field of Section~\ref{sec:3.1} smoothly across the
axis, we factor out the linear radial vanishing of its azimuthal component. We
introduce a fixed amplitude normalization \NSi{NS-F-P025-017}{\mathsf C>1} and a
smooth scalar \NSi{NS-F-P025-018}{\phi}, and write
\NSFormula{NS-F-P025-019}{display}{4.3}
\begin{equation}
 u_\theta^{(0)}=q^{-A}E,\qquad
 u_z^{(0)}=q^{-A}U,\qquad
 ru_r^{(0)}=V_0,\qquad
 p^{(0)}=q^{-2A}\Pi,\qquad
 E=\mathsf C^{-1}\sqrt{2X}\,\phi.
\tag{4.3}\label{eq:4.3}
\end{equation}
The normalization \NSi{NS-F-P025-020}{\mathsf C} is chosen in the proof of
Theorem~\ref{thm:4.6}. As in Section~\ref{sec:3.1},
\NSi{NS-F-P025-021}{E,U,V_0,\Pi} are functions of
\NSi{NS-F-P025-022}{(X,\eta)}. The factor
\NSi{NS-F-P025-023}{\sqrt{2X}} accounts for the vanishing of the azimuthal
velocity at the axis; smoothness is imposed on
\NSi{NS-F-P025-024}{\phi}. For the profiles constructed here, regularity at
the axis in the sense of Definition~\ref{def:3.2} follows from the conditions
\NSFormula{NS-F-P025-025}{display}{4.4}
\begin{equation}
 E=\sqrt{2X}\,F,\qquad V_0=Xv_0,\qquad
 F=\phi/\mathsf C,\qquad
 U,v_0,\Pi\in C^\infty([0,X_c]\times[-1,1])
\tag{4.4}\label{eq:4.4}
\end{equation}
for some \NSi{NS-F-P025-026}{X_c>0}. Smoothness on this closed rectangle
means that all mixed \NSi{NS-F-P025-027}{X,\eta} derivatives have continuous
one-sided extensions to its boundary. To see the Cartesian smoothness, write
\NSi{NS-F-P025-028}{x_1=r\cos\theta,\ x_2=r\sin\theta}, so that
\NSi{NS-F-P025-029}{X=(x_1^2+x_2^2)/(2q)}. Then
\NSFormula{NS-F-P025-030}{display}{4.5}
\begin{equation}
\begin{aligned}
 u_1^{(0)}&=\frac{v_0}{2q}x_1-q^{-A-1/2}Fx_2,
 &u_2^{(0)}&=\frac{v_0}{2q}x_2+q^{-A-1/2}Fx_1,\\
 u_3^{(0)}&=q^{-A}U,
 &p^{(0)}&=q^{-2A}\Pi.
\end{aligned}
\tag{4.5}\label{eq:4.5}
\end{equation}
Here all profiles are evaluated at \NSi{NS-F-P025-031}{(X,\eta)}, and
\NSi{NS-F-P025-032}{q,\eta} are smooth functions of
\NSi{NS-F-P025-033}{(t,z)} with \NSi{NS-F-P025-034}{q>0} for
\NSi{NS-F-P025-035}{t<1}. Thus the coefficients in \eqref{eq:4.5} are smooth
functions of \NSi{NS-F-P025-036}{(x_1,x_2,z,t)}. In particular,
\NSi{NS-F-P025-037}{E} itself contains a \NSi{NS-F-P025-038}{\sqrt X} factor;
smoothness at \NSi{NS-F-P025-039}{X=0} is imposed on
\NSi{NS-F-P025-040}{F}.

We choose \NSi{NS-F-P025-041}{E,U}; incompressibility determines
\NSi{NS-F-P025-042}{V_0}, and the leading radial momentum balance determines
\NSi{NS-F-P025-043}{\Pi} once its axis value is fixed. For a smooth radial
scalar, recall the radial average from the notation paragraph at the end of
Section~\ref{sec:3}; we specify its smooth extension to the axis by
\NSFormula{NS-F-P025-044}{display}{4.6}
\begin{equation}
 \mathcal A_X(f)(X,\eta)=\frac1X\int_0^X f(x,\eta)\,dx,
 \qquad
 \mathcal A_X(f)(0,\eta)=f(0,\eta).
\tag{4.6}\label{eq:4.6}
\end{equation}

\NSPage{026}%
The incompressibility and centrifugal-pressure balances stated in
Section~\ref{sec:3.1} give the following explicit profile identities:
\NSFormula{NS-F-P026-001}{display}{4.7}
\begin{equation}
 V_0=\frac{X}{L}\bigl(2\eta U-2D\eta\mathcal A_X(U)
      -d\partial_\eta\mathcal A_X(U)\bigr),
 \qquad
 \Pi_X=\frac{E^2}{2X}.
\tag{4.7}\label{eq:4.7}
\end{equation}
To see the first identity, use \NSi{NS-F-P026-002}{A+D=1} in
\NSi{NS-F-P026-003}{\partial_s(ru_r^{(0)})+\partial_z u_z^{(0)}=0}. It gives
\NSi{NS-F-P026-004}{(V_0)_X=L^{-1}(2A\eta U-dU_\eta+2\eta XU_X)};
integration with zero axis value gives the displayed formula. For the second
identity, the coefficients of \NSi{NS-F-P026-005}{\partial_rp^{(0)}} and
\NSi{NS-F-P026-006}{(u_\theta^{(0)})^2/r} are respectively
\NSi{NS-F-P026-007}{\sqrt{2X}\Pi_X} and
\NSi{NS-F-P026-008}{E^2/\sqrt{2X}}, with common power
\NSi{NS-F-P026-009}{q^{-2A-1/2}}. The remaining radial terms, and axial
viscosity relative to radial viscosity in the tangential equations, carry an
additional factor at least \NSi{NS-F-P026-010}{q^{2h}}. This viscosity
comparison explains the powers \NSi{NS-F-P026-011}{q^{2nh}} used in the
background corrections of Section~\ref{sec:3}. Here these are comparisons at
fixed profile coordinates; physical derivative bounds will be proved
separately for the completed fields. The higher-order terms just identified
are retained in the full residual and treated in Section~\ref{sec:5}.

With incompressibility and the leading radial balance imposed, it remains to
express the tangential residual as a radial stress divergence, as required in
Section~\ref{sec:3}. We first integrate its inviscid terms, denoting the
resulting scalar profiles by \NSi{NS-F-P026-012}{Q_s,N_s}, and then add the
contribution from radial viscosity. The definitions below depend only on
\NSi{NS-F-P026-013}{E,U,\Pi}; Proposition~\ref{prop:4.2} verifies the resulting
stress identity.

For profiles satisfying \eqref{eq:4.7}, with
\NSi{NS-F-P026-014}{\phi,U,\Pi} smooth at the axis and
\NSi{NS-F-P026-015}{\phi>0}, we set
\NSFormula{NS-F-P026-016}{display}{4.8}
\begin{equation}
\begin{gathered}
 H=\sqrt{2X}E,\qquad F=\frac{E}{\sqrt{2X}},\qquad l=\mathcal D_X\log H,\\
 W=1-2D\eta\mathcal A_X(U)-d\partial_\eta\mathcal A_X(U),
 \qquad H_c=D\eta+dU.
\end{gathered}
\tag{4.8}\label{eq:4.8}
\end{equation}
Here \NSi{NS-F-P026-017}{H} is the profile of angular momentum
\NSi{NS-F-P026-018}{ru_\theta^{(0)}},
\NSi{NS-F-P026-019}{F} removes the factor
\NSi{NS-F-P026-020}{\sqrt{2X}} from \NSi{NS-F-P026-021}{E}, and
\NSi{NS-F-P026-022}{l} measures the logarithmic radial slope of
\NSi{NS-F-P026-023}{H}. The coefficients
\NSi{NS-F-P026-024}{W,H_c} collect the transport terms in \eqref{eq:4.14}.
We define \NSi{NS-F-P026-025}{Q_s,N_s} by the following radial equations,
selecting the solutions that extend smoothly to \NSi{NS-F-P026-026}{X=0}:
\NSFormula{NS-F-P026-027}{display}{4.9}
\begin{equation}
\begin{aligned}
 \mathcal D_XQ_s+(1+l)Q_s&=S_q,& \mathcal D_XN_s+N_s&=S_n,\\
 S_q&=-Wl-h(1-2\eta U)-H_c(\log E)_\eta,\\
 S_n&=-W\mathcal D_XU-A(1-2\eta U)U-H_cU_\eta
      -d\Pi_\eta+4A\eta\Pi+2\eta \mathcal D_X\Pi.
\end{aligned}
\tag{4.9}\label{eq:4.9}
\end{equation}
The integrating factors are \NSi{NS-F-P026-028}{XH} and
\NSi{NS-F-P026-029}{X}, respectively. Thus the integration constants, which
may a priori depend on \NSi{NS-F-P026-030}{\eta}, enter as
\NSFormula{NS-F-P026-031}{display}{4.10}
\begin{equation}
 Q_s=\frac{C_Q(\eta)+\int_0^X H(x,\eta)S_q(x,\eta)\,dx}{XH(X,\eta)},
 \qquad
 N_s=\frac{C_N(\eta)+\int_0^X S_n(x,\eta)\,dx}{X}.
\tag{4.10}\label{eq:4.10}
\end{equation}
The required smooth extensions select
\NSi{NS-F-P026-032}{C_Q=C_N=0}. Indeed,
\NSi{NS-F-P026-033}{H=2X\phi/\mathsf C} near the axis, a nonzero constant would
produce an \NSi{NS-F-P026-034}{X^{-2}} singularity in
\NSi{NS-F-P026-035}{Q_s}, or an \NSi{NS-F-P026-036}{X^{-1}} singularity in
\NSi{NS-F-P026-037}{N_s}. With both constants zero, changing variables
\NSi{NS-F-P026-038}{x=sX} gives
\NSFormula{NS-F-P026-039}{display}{none}
\[
 Q_s=\frac{1}{F(X,\eta)}\int_0^1 sF(sX,\eta)S_q(sX,\eta)\,ds,
 \qquad
 N_s=\int_0^1 S_n(sX,\eta)\,ds.
\]
The sources are smooth at \NSi{NS-F-P026-040}{X=0}, since
\NSi{NS-F-P026-041}{l=1+XF_X/F} and
\NSi{NS-F-P026-042}{(\log E)_\eta=F_\eta/F}. These formulas give the smooth
extensions \NSi{NS-F-P026-043}{Q_s(0,\eta)=S_q(0,\eta)/2} and
\NSi{NS-F-P026-044}{N_s(0,\eta)=S_n(0,\eta)}. We define the two
\NSPage{027}%
stress coefficients in the order \NSi{NS-F-P027-001}{(r\theta,rz)} by
\NSFormula{NS-F-P027-002}{display}{4.11}
\begin{equation}
\begin{gathered}
 \mathcal T_0=F(p_s-\mathfrak{s}),
 \qquad
 p_s=\left(\frac{XQ_s}{L},\frac{XN_s}{LE}\right),\\
 \mathfrak{s}=(a,-b_s),
 \qquad
 a=1-2\mathcal D_X\log E=2-2l,
 \qquad
 b_s=\frac{2\mathcal D_XU}{E}.
\end{gathered}
\tag{4.11}\label{eq:4.11}
\end{equation}
Thus \NSi{NS-F-P027-003}{p_s} is a two-component vector of integrated inviscid
contributions, distinct from the pressure \NSi{NS-F-P027-004}{p^{(0)}}. The
vector \NSi{NS-F-P027-005}{\mathfrak{s}=(a,-b_s)} records the cylindrical
radial shear, with its normalization specified by
\NSFormula{NS-F-P027-006}{display}{none}
\[
 \left(\partial_ru_\theta^{(0)}-\frac{u_\theta^{(0)}}{r},
       \partial_ru_z^{(0)}\right)
 =-q^{-A-1/2}F\mathfrak{s}.
\]
In \eqref{eq:4.11}, \NSi{NS-F-P027-007}{Fp_s} is the inviscid contribution
and \NSi{NS-F-P027-008}{-F\mathfrak{s}} is the radial-viscosity contribution,
as the calculation below shows. In the notation of Section~\ref{sec:3}, the
physical stress representation to be proved is
\NSi{NS-F-P027-009}{T=q^{-A-1/2}\mathcal T_0}. Its two components are the
required \NSi{NS-F-P027-010}{r\theta} and \NSi{NS-F-P027-011}{rz} entries of the wave
covariance, the tensor of averaged quadratic velocity products constructed in
Proposition~\ref{prop:7.5}. To verify the stress identity used in
Section~\ref{sec:3}, we must separate the tangential terms at the leading
order from axial viscosity. Define the \emph{leading tangential residuals} by
\NSFormula{NS-F-P027-012}{display}{4.12}
\begin{equation}
 R_j^{(0)}=\mathscr R(u^{(0)},p^{(0)})\mathbin{\cdot}e_j
            +\partial_z^2u_j^{(0)},
 \qquad j\in\{\theta,z\},
\tag{4.12}\label{eq:4.12}
\end{equation}
where \NSi{NS-F-P027-013}{e_\theta,e_z} are the cylindrical unit vectors and
\NSi{NS-F-P027-014}{\mathscr R} is the full momentum residual defined in the
introduction. Adding \NSi{NS-F-P027-015}{\partial_z^2u_j^{(0)}} removes axial
viscosity from that residual; radial viscosity is retained.

\begin{proposition}\label{prop:4.2}
The residuals \NSi{NS-F-P027-016}{R_\theta^{(0)},R_z^{(0)}} in
\eqref{eq:4.12} are minus the cylindrical divergences
\NSi{NS-F-P027-017}{\partial_r+2/r} and
\NSi{NS-F-P027-018}{\partial_r+1/r} of
\NSi{NS-F-P027-019}{q^{-A-1/2}\mathcal T_0}. In particular the equations for
vanishing leading residual stress are
\NSFormula{NS-F-P027-020}{display}{4.13}
\begin{equation}
 -2L\frac{X\phi_{XX}+2\phi_X}{\phi}=S_q,
 \qquad
 -2L(XU_{XX}+U_X)=S_n.
\tag{4.13}\label{eq:4.13}
\end{equation}
\end{proposition}

\begin{proof}
The angular equation is simplest in angular momentum
\NSi{NS-F-P027-021}{ru_\theta^{(0)}=q^{-h}H}. Equations \eqref{eq:4.2} and
\eqref{eq:4.7} give
\NSFormula{NS-F-P027-022}{display}{4.14}
\begin{equation}
\begin{split}
 &(\partial_t+u_r^{(0)}\partial_r+u_z^{(0)}\partial_z)(q^{-h}H)\\
 &\qquad=\frac{q^{-h-1}}{L}
   \{W\mathcal D_XH+H_cH_\eta+h(1-2\eta U)H\}
   =-\frac{q^{-h-1}}{L}HS_q.
\end{split}
\tag{4.14}\label{eq:4.14}
\end{equation}
The axial material derivative plus the pressure gradient is similarly
\NSi{NS-F-P027-023}{-q^{-A-1}S_n/L}. The integrating factors for the two
radial divergences are \NSi{NS-F-P027-024}{r^2} and
\NSi{NS-F-P027-025}{r}. Integrating the negative inviscid residual from the
axis therefore produces the coefficients \NSi{NS-F-P027-026}{FXQ_s/L} and
\NSi{NS-F-P027-027}{FXN_s/(LE)}. Radial viscosity contributes
\NSFormula{NS-F-P027-028}{display}{none}
\[
 \partial_ru_\theta^{(0)}-u_\theta^{(0)}/r=-q^{-A-1/2}Fa,
 \qquad
 \partial_ru_z^{(0)}=q^{-A-1/2}Fb_s.
\]
Adding these viscous contributions to the radially integrated inviscid
residual proves \eqref{eq:4.11}, including its sign. If that stress is zero
then \NSi{NS-F-P027-029}{XQ_s/L=a,\ N_s/L=-2U_X}; inserting these relations
in \eqref{eq:4.9} proves \eqref{eq:4.13}. At the axis
\NSi{NS-F-P027-030}{l=1+X\phi_X/\phi} and
\NSi{NS-F-P027-031}{(\log E)_\eta=\phi_\eta/\phi} are smooth, so
\NSi{NS-F-P027-032}{Q_s,N_s} are regular at the axis in the scalar-profile
sense of Definition~\ref{def:3.2}.
\end{proof}

\NSPage{028}%
\subsection{The cumulative radial integrals}\label{sec:4.2}

Note that the stress profile \NSi{NS-F-P028-001}{\mathcal T_0} in
\eqref{eq:4.11} depends on the profiles \NSi{NS-F-P028-002}{E,U} at every
smaller radius through the integrals in \eqref{eq:4.10}. To join an inner
profile to a prescribed outer one, as in Section~\ref{sec:3}, we must therefore
match the accumulated integrals as well as the profile values. Five integrals
suffice to preserve the outer pressure, radial velocity, and stress
(Lemma~\ref{lem:4.4}). Integrating \eqref{eq:4.9} expresses
\NSi{NS-F-P028-003}{Q_s,N_s} in terms of these quantities:
\NSFormula{NS-F-P028-004}{display}{4.15}
\begin{equation}
\begin{gathered}
 M=\int_0^X U\,dx,\qquad
 I=\int_0^X H\,dx,\qquad
 J=\int_0^X UH\,dx,\\
 S=\int_0^X(U^2-E^2/2)\,dx,\qquad
 C_p=\int_0^X\frac{E^2}{2x}\,dx,\qquad
 \Pi=\Pi(0,\eta)+C_p.
\end{gathered}
\tag{4.15}\label{eq:4.15}
\end{equation}
Each integral is a function of \NSi{NS-F-P028-005}{(X,\eta)}; at a fixed
matching radius it is a scalar function of \NSi{NS-F-P028-006}{\eta}. Here
\NSi{NS-F-P028-007}{C_p} is the pressure increment from the axis to
\NSi{NS-F-P028-008}{X}, and \NSi{NS-F-P028-009}{\Pi(0,\eta)} is the
prescribed value of the pressure at \NSi{NS-F-P028-010}{X=0}.

\begin{lemma}\label{lem:4.3}
For the solutions of \eqref{eq:4.9} with
\NSi{NS-F-P028-011}{C_Q=C_N=0},
\NSFormula{NS-F-P028-012}{display}{4.16}
\begin{equation}
\begin{aligned}
 Q_s&=-W+\frac{(1-h)I-D\eta I_\eta-dJ_\eta+2(h-D)\eta J}{XH},\\
 N_s&=-WU+\frac{D(M-\eta M_\eta)+4h\eta S-dS_\eta}{X}
             +4A\eta\Pi-d\Pi_\eta.
\end{aligned}
\tag{4.16}\label{eq:4.16}
\end{equation}
\end{lemma}

\begin{proof}
Use \NSi{NS-F-P028-013}{\partial_X(XW)=1-2D\eta U-dU_\eta} to integrate the
terms \NSi{NS-F-P028-014}{-W\mathcal D_XH} and
\NSi{NS-F-P028-015}{-W\mathcal D_XU}. The angular terms
\NSi{NS-F-P028-016}{-dU_\eta H-dUH_\eta} combine to
\NSi{NS-F-P028-017}{-d(UH)_\eta}. For the pressure contribution, integration
by parts in \NSi{NS-F-P028-018}{\Pi_X=E^2/(2X)} gives
\NSi{NS-F-P028-019}{\int_0^X\Pi\,dx=X\Pi-\int_0^XE^2/2\,dx}.
Consequently the source integrals are
\NSFormula{NS-F-P028-020}{display}{none}
\[
\begin{aligned}
 \int_0^XHS_q\,dx
   &=-XWH+(1-h)I-D\eta I_\eta-dJ_\eta+2(h-D)\eta J,\\
 \int_0^XS_n\,dx
   &=-XWU+D(M-\eta M_\eta)+4h\eta S-dS_\eta
      +X(4A\eta\Pi-d\Pi_\eta).
\end{aligned}
\]
Division by \NSi{NS-F-P028-021}{XH} and \NSi{NS-F-P028-022}{X} gives the result.
\end{proof}

When two profiles agree beyond a joining radius, matching the five cumulative
integrals there also preserves the quantities obtained by radial integration.
Their common pressure value at the axis is held fixed. Small changes in the
profiles, integrals, and their parameter derivatives also give small changes
in \NSi{NS-F-P028-023}{Q_s,N_s,p_s}, without requiring radial derivatives of
the profiles to be close. We use the exact identity to join profiles and the
estimate to alter their shear.

\begin{lemma}\label{lem:4.4}
Fix \NSi{NS-F-P028-024}{h\in(0,1/2)}. For
\NSi{NS-F-P028-025}{i=1,2}, let \NSi{NS-F-P028-026}{(U_i,E_i)} be smooth on
\NSi{NS-F-P028-027}{(0,\infty)\times[-1,1]}, with
\NSi{NS-F-P028-028}{E_i>0}, and suppose the five integrals in
\eqref{eq:4.15} define smooth functions there, including one-sided
\NSi{NS-F-P028-029}{\eta}-derivatives at
\NSi{NS-F-P028-030}{\eta=\pm1}. Write
\NSFormula{NS-F-P028-031}{display}{none}
\[
 H_i=\sqrt{2X}E_i,\qquad
 \mathfrak m_i=(M_i,I_i,J_i,S_i,C_{p,i}).
\]
Fix one smooth function \NSi{NS-F-P028-032}{\Pi_{\mathrm{ax}}(\eta)} on
\NSi{NS-F-P028-033}{[-1,1]} and set
\NSFormula{NS-F-P028-034}{display}{none}
\[
 \Pi_i(X,\eta)=\Pi_{\mathrm{ax}}(\eta)+C_{p,i}(X,\eta),
 \qquad
 \Pi_i(0,\eta)=\Pi_{\mathrm{ax}}(\eta),
 \qquad i=1,2.
\]
Thus the pressures have the same integration constant at the axis; the
velocity profiles need not agree near the axis. Define
\NSi{NS-F-P028-035}{V_{0,i}} by \eqref{eq:4.7},
\NSi{NS-F-P028-036}{Q_{s,i},N_{s,i}} by \eqref{eq:4.16}, and
\NSi{NS-F-P028-037}{p_{s,i},a_i,b_{s,i},\mathcal T_{0,i}} by
\eqref{eq:4.11}. We use the explicit formulas \eqref{eq:4.16}, previously
derived with \NSi{NS-F-P028-038}{C_Q=C_N=0}. For any of these quantities
\NSi{NS-F-P028-039}{f}, let
\NSi{NS-F-P028-040}{\Delta f=f_2-f_1}. Then the following two conclusions
hold.
\NSPage{029}%
\begin{enumerate}[label=\textup{(\roman*)},leftmargin=2.5em]
\item Suppose that for some \NSi{NS-F-P029-001}{X_h>0},
\NSFormula{NS-F-P029-002}{display}{none}
\[
 (U_1,E_1)=(U_2,E_2)\quad\text{for }X\geq X_h,
 \qquad
 \Delta\mathfrak m(X_h,\eta)=0\quad(-1\leq\eta\leq1).
\]
Then \NSi{NS-F-P029-003}{\mathfrak m_1=\mathfrak m_2}, and the corresponding
quantities
\NSi{NS-F-P029-004}{\Pi_i,V_{0,i},Q_{s,i},N_{s,i},p_{s,i},a_i,b_{s,i},
\mathcal T_{0,i}} agree for all \NSi{NS-F-P029-005}{X\geq X_h} and
\NSi{NS-F-P029-006}{\eta\in[-1,1]}.

\item Fix \NSi{NS-F-P029-007}{0<X_0<X_1<\infty}, an integer
\NSi{NS-F-P029-008}{k\geq0}, and \NSi{NS-F-P029-009}{e_{\min}>0}. On
\NSi{NS-F-P029-010}{\mathcal R=[X_0,X_1]\times[-1,1]}, put
\NSFormula{NS-F-P029-011}{display}{none}
\[
 \|f\|_j:=\max_{0\leq r\leq j}
           \sup_{(X,\eta)\in\mathcal R}|\partial_\eta^rf(X,\eta)|,
 \qquad j\geq0,
\]
using the maximum of the component norms for tuples. If
\NSi{NS-F-P029-012}{E_i\geq e_{\min}} on
\NSi{NS-F-P029-013}{\mathcal R}, then
\NSFormula{NS-F-P029-014}{display}{4.17}
\begin{equation}
 \|\Delta(Q_s,N_s,p_s)\|_k
 \leq C_k\bigl(\|\Delta(U,E)\|_{k+1}
                 +\|\Delta\mathfrak m\|_{k+1}\bigr).
\tag{4.17}\label{eq:4.17}
\end{equation}
The constant \NSi{NS-F-P029-015}{C_k} depends only on
\NSi{NS-F-P029-016}{k,h,X_0,X_1,e_{\min}}, the
\NSi{NS-F-P029-017}{C^{k+1}}-norm of
\NSi{NS-F-P029-018}{\Pi_{\mathrm{ax}}}, and a common bound for
\NSi{NS-F-P029-019}{\|(U_i,E_i,\mathfrak m_i)\|_{k+1}},
\NSi{NS-F-P029-020}{i=1,2}. No radial derivative of a difference occurs in
this estimate.
\end{enumerate}
\end{lemma}

\begin{proof}
By the definitions of the five cumulative integrals,
\NSFormula{NS-F-P029-021}{display}{4.18}
\begin{equation}
\begin{aligned}
 \Delta M&=\int_0^X(U_2-U_1)\,dx,\\
 \Delta I&=\int_0^X(H_2-H_1)\,dx,\\
 \Delta J&=\int_0^X(U_2H_2-U_1H_1)\,dx,\\
 \Delta S&=\int_0^X\bigl(U_2^2-U_1^2-(E_2^2-E_1^2)/2\bigr)\,dx,\\
 \Delta C_p&=\int_0^X\frac{E_2^2-E_1^2}{2x}\,dx.
\end{aligned}
\tag{4.18}\label{eq:4.18}
\end{equation}
The quantities on the left are evaluated at
\NSi{NS-F-P029-022}{(X,\eta)}, and each integrand on the right at
\NSi{NS-F-P029-023}{(x,\eta)}. For part~\textup{(i)}, agreement of
\NSi{NS-F-P029-024}{U_i,E_i} on \NSi{NS-F-P029-025}{X\geq X_h} makes each
integrand in \eqref{eq:4.18} zero there. Consequently
\NSFormula{NS-F-P029-026}{display}{none}
\[
 \Delta\mathfrak m(X,\eta)=\Delta\mathfrak m(X_h,\eta)=0
 \qquad(X\geq X_h,\ -1\leq\eta\leq1).
\]
This is an identity of functions of \NSi{NS-F-P029-027}{\eta}, so their
\NSi{NS-F-P029-028}{\eta}-derivatives agree as well. The common pressure datum
gives \NSi{NS-F-P029-029}{\Delta\Pi=\Delta C_p=0}. Equation \eqref{eq:4.7}
then gives \NSi{NS-F-P029-030}{\Delta V_0=0}, and \eqref{eq:4.16} gives
\NSi{NS-F-P029-031}{\Delta Q_s=\Delta N_s=0}. The profiles agree on a radial
interval, so their radial derivatives also agree. The definitions in
\eqref{eq:4.11} therefore give equality of
\NSi{NS-F-P029-032}{p_s,a,b_s,\mathcal T_0}.

For part~\textup{(ii)}, the quantities entering \eqref{eq:4.16} can be written
as
\NSFormula{NS-F-P029-033}{display}{none}
\[
 H_i=\sqrt{2X}E_i,
 \qquad
 W_i=1-\frac{2D\eta M_i+d\partial_\eta M_i}{X},
 \qquad
 \Pi_i=\Pi_{\mathrm{ax}}+C_{p,i}.
\]
It follows that
\NSFormula{NS-F-P029-034}{display}{none}
\[
 \Delta H=\sqrt{2X}\Delta E,
 \qquad
 \Delta W=-\frac{2D\eta\Delta M+d\partial_\eta\Delta M}{X},
 \qquad
 \Delta\Pi=\Delta C_p.
\]
The last identity is the use of the common axis pressure datum. On
\NSi{NS-F-P029-035}{\mathcal R}, the denominators satisfy
\NSi{NS-F-P029-036}{X\geq X_0},
\NSi{NS-F-P029-037}{L\geq1-2h>0}, and
\NSi{NS-F-P029-038}{H_i\geq\sqrt{2X_0}e_{\min}}. For example,
\NSFormula{NS-F-P029-039}{display}{none}
\[
 \Delta(H^{-1})=-\frac{\Delta H}{H_1H_2},
 \qquad
 \Delta(E^{-1})=-\frac{\Delta E}{E_1E_2}.
\]
Apply the product rule and these reciprocal identities to \eqref{eq:4.16},
and then to the formula for \NSi{NS-F-P029-040}{p_s} in \eqref{eq:4.11}.
After at most \NSi{NS-F-P029-041}{k}
\NSi{NS-F-P029-042}{\eta}-derivatives, each term contains a difference of
\NSi{NS-F-P029-043}{U,E}, or of
\NSPage{030}%
a component of \NSi{NS-F-P030-001}{\mathfrak m}, through order at most
\NSi{NS-F-P030-002}{k+1}. The remaining factors are controlled by the stated
common bounds. This proves \eqref{eq:4.17} and the dependence of its constant.
These formulas contain no radial derivatives of the inputs.
\end{proof}

Part~\textup{(ii)} applies to the radial modulation in
Proposition~\ref{prop:C.2}. The profiles \NSi{NS-F-P030-003}{U,E}, the five
integrals \NSi{NS-F-P030-004}{\mathfrak m}, and their
\NSi{NS-F-P030-005}{\eta}-derivatives remain close, even though the radial
shears \NSi{NS-F-P030-006}{a,b_s} change substantially. The estimate controls
\NSi{NS-F-P030-007}{p_s}; the full stress
\NSi{NS-F-P030-008}{\mathcal T_0=F(p_s-\mathfrak{s})} also depends on those
shears.

To use this comparison at a large joining radius
\NSi{NS-F-P030-009}{\rho}, we set
\NSi{NS-F-P030-010}{\widehat U(x,\eta)=U(\rho x,\eta)} and
\NSi{NS-F-P030-011}{\widehat E(x,\eta)=E(\rho x,\eta)}, and define the
remaining hatted quantities by
\NSFormula{NS-F-P030-012}{display}{4.19}
\begin{equation}
 X=\rho x,\qquad H=\sqrt\rho\,\widehat H,\qquad
 (M,I,J,S,C_p)
   =(\rho\widehat M,\rho^{3/2}\widehat I,\rho^{3/2}\widehat J,
     \rho\widehat S,\widehat C_p).
\tag{4.19}\label{eq:4.19}
\end{equation}
Here unhatted functions are evaluated at
\NSi{NS-F-P030-013}{(\rho x,\eta)}, and hatted functions at
\NSi{NS-F-P030-014}{(x,\eta)}. We put
\NSFormula{NS-F-P030-015}{display}{none}
\[
\begin{gathered}
 \widehat{\mathfrak m}
   =(\widehat M,\widehat I,\widehat J,\widehat S,\widehat C_p),\\
 (\widehat Q_s,\widehat N_s)(x,\eta)=(Q_s,N_s)(\rho x,\eta),\\
 \widehat p_s(x,\eta)=\rho^{-1}p_s(\rho x,\eta).
\end{gathered}
\]
Substitution into \eqref{eq:4.16} and \eqref{eq:4.11} gives the same formulas
for \NSi{NS-F-P030-016}{\widehat Q_s,\widehat N_s,\widehat p_s} with
\NSi{NS-F-P030-017}{X} replaced by \NSi{NS-F-P030-018}{x}. Thus on a fixed
rectangle \NSi{NS-F-P030-019}{[x_0,x_1]\times[-1,1]}, with
\NSi{NS-F-P030-020}{0<x_0<x_1},
\NSFormula{NS-F-P030-021}{display}{none}
\[
 \|\Delta(\widehat Q_s,\widehat N_s,\widehat p_s)\|_k
 \leq C_k\bigl(\|\Delta(\widehat U,\widehat E)\|_{k+1}
                 +\|\Delta\widehat{\mathfrak m}\|_{k+1}\bigr),
\]
where the norms use this fixed \NSi{NS-F-P030-022}{x}-rectangle. The constant
is independent of \NSi{NS-F-P030-023}{\rho} provided the positive lower bound
for \NSi{NS-F-P030-024}{\widehat E_i}, the bounds for
\NSi{NS-F-P030-025}{\|(\widehat U_i,\widehat E_i,
\widehat{\mathfrak m}_i)\|_{k+1}}, and the
\NSi{NS-F-P030-026}{C^{k+1}}-bound for the common
\NSi{NS-F-P030-027}{\Pi_{\mathrm{ax}}} are independent of
\NSi{NS-F-P030-028}{\rho}, with \NSi{NS-F-P030-029}{h} fixed.

The moment corrections used below add finite linear combinations of smooth
functions supported on short radial intervals. Their changes to
\NSi{NS-F-P030-030}{\mathfrak m} are linear or quadratic in the coefficients.
The invertibility of the moment matrices and the solution of these finite
systems are proved in Lemmas~\ref{lem:A.1} and~\ref{lem:A.2}; at each
application we verify the required bounds for the moment discrepancies and
inverses.

Moment conditions at infinity serve a separate purpose. The quantities
\NSi{NS-F-P030-031}{M,J,S} control, respectively, the radially integrated
axial momentum, the axial transport of angular momentum, and the axial
momentum flux including pressure. The angular moment
\NSi{NS-F-P030-032}{I} is made convergent by subtracting the prescribed
exterior power \NSi{NS-F-P030-033}{H_{\mathrm{pow}}} before integration, as
in \eqref{eq:4.28}. The resulting total identities make the two weighted
integrals of the tangential residuals vanish. Without them, a zero exterior
residual could leave stresses proportional to \NSi{NS-F-P030-034}{r^{-2}}
and \NSi{NS-F-P030-035}{r^{-1}} in the angular and axial components: the
integrated inner residual may remain nonzero. Lemma~\ref{lem:A.8} uses the
total moment identities to remove these tails and prove that the stress
vanishes in the exterior.

\subsection{The cone imposed on the leading stress}\label{sec:4.3}

Since the annular stress \NSi{NS-F-P030-036}{\mathcal T_0} must be produced
by the chosen waves with real amplitudes, it must lie in the positive span of
their covariance directions (Proposition~\ref{prop:7.5}). We now express the
required cone condition in the profile variables; Lemma~\ref{lem:4.5} gives
the equivalent inequalities. At each \NSi{NS-F-P030-037}{(X,\eta)}, the
relevant profile data are the radial shear
\NSi{NS-F-P030-038}{(a,-b_s)} and the integrated inviscid contribution
\NSi{NS-F-P030-039}{p_s=(p_{s,1},p_{s,2})} from \eqref{eq:4.11}. Where
\NSi{NS-F-P030-040}{a>0}, we set
\NSFormula{NS-F-P030-041}{display}{4.20}
\begin{equation}
 t_s=-\frac{b_s}{a},\qquad
 v_s=a(1+t_s^2),\qquad
 P_c=p_{s,1}+t_sp_{s,2},\qquad
 J_c=p_{s,2}-t_sp_{s,1}.
\tag{4.20}\label{eq:4.20}
\end{equation}
